\chapter{MSX and Amiga Compatibility Checklist}

\section{BX-8, in place of an MSX}

\begin{itemize}[leftmargin=1.5em]
  \item Eight-bit data bus shared by the hart's data port, the video display processor, and the sound generator.
  \item Slot and cartridge signals documented against the MSX reference you follow.
  \item Video modes you claim: register map, timing, and a captured frame compared to that reference.
  \item Sound: tone, noise, and envelope registers, and a captured waveform.
  \item Keyboard and joystick ports in the documented map.
  \item Video refresh can take the bus during a vector load; the picture remains stable; the load's bytes match a quiet-bus run.
  \item The hart is RV32I with SVM-8. Guest binaries of the original processor run only through software you have written for that purpose, and that software is named on the card.
\end{itemize}

\section{BX-16, in place of an Amiga}

\begin{itemize}[leftmargin=1.5em]
  \item Sixteen-bit shared bus with hart, blitter-class master, video, and audio.
  \item A reserved-slot or priority rule for audio, with a measured worst-case wait.
  \item Display fetch can be driven from a programmed list.
  \item Blitter-class master, if present, is a peripheral, not a second instruction set.
  \item A vector fill of halfwords matches a scalar reference while audio slots are taken.
  \item The hart is RV32I with SVM-16. Xsvm is off on the compatibility image.
\end{itemize}

\section{BX-32 and BX-64}

These images are compatible with the RISC-V strings on their profile cards, and with the travelling tests. They are not MSX or Amiga replacements. Their checklist is the RVV checklist plus the memory map you publish.
